\documentclass[11pt,a4paper,twocolumn]{article}

\usepackage[utf8]{inputenc}
\usepackage[T1]{fontenc}
\usepackage[a4paper,margin=2.2cm]{geometry}
\usepackage{amsmath,amssymb}
\usepackage{booktabs}
\usepackage{graphicx}
\usepackage{enumitem}
\usepackage{array}
\usepackage{dblfloatfix}
\usepackage{microtype}
\usepackage[numbers,sort&compress]{natbib}
\usepackage[hidelinks]{hyperref}

\setlist{nosep,leftmargin=1.2em}
\newcommand{\vect}[1]{\boldsymbol{#1}}

\begin{document}

% Full-width front matter, then two columns.
\twocolumn[%
\begin{center}
  {\LARGE\bfseries Multi-Objective Optimization for Large-Scale\\[0.25em]
   Neural Solver Deployment for Combinatorial Optimization\par}
  \vspace{1.1em}
  {\large Yuxuan LI\par}
  \vspace{0.25em}
  {\normalsize\textit{Supervisor:} Ruibin BAI\par}
\end{center}
\vspace{0.9em}
\begin{center}
\begin{minipage}{0.94\textwidth}
\small\noindent\textbf{Abstract.} Neural combinatorial optimization (NCO) reports
solution quality at a single operating point, leaving system-level questions
unanswered: where inference time goes, how much of it is recoverable, and whether
deployment is worth its cost. This proposal recasts ``large-scale'' in NCO as a
\emph{deployment-scale} property --- instances served, concurrency, and sustained
operating duration --- rather than instance size alone. Work proceeds in four
layers: (A)~characterization of where inference cost is spent and how hotspots
migrate with scale; (B)~quantification of recoverable headroom from compilation,
batching, and precision; (C)~a multi-objective formulation yielding Pareto-optimal
operating points and a deployment decision rule under a quality SLA; and
(D)~a generalization study across problem family and execution backend. A central
methodological contribution is a \emph{hardware calibration anchor} that converts
cross-backend latency ratios into a normalized efficiency metric, separating three
findings a naive comparison conflates: the hardware is slower, the backend is
inefficient, or the workload is bound by neither and dominated by launch and
synchronization overhead.
\end{minipage}
\end{center}
\vspace{1.2em}
]

%======================================================================
\section{Introduction}
%======================================================================

Neural combinatorial optimization (NCO) has matured into a rich landscape of
learned solvers: autoregressive construction \citep{kool2019attention,kwon2020pomo},
heatmap prediction with search \citep{joshi2019gnn,sun2023difusco}, improvement and
local search \citep{ma2021dact,hottung2022eas}, and decomposition for large
instances \citep{duan2023htsp,ye2024glop}. Reporting conventions have not kept
pace. Papers typically give one (quality, runtime) point per instance size, on one
accelerator, with runtime as an undifferentiated wall clock. Three questions that
decide whether such a solver can actually be deployed are therefore unanswered.

\emph{Where does the time go?} Autoregressive decoders execute a serial loop of
$O(n)$ steps, each dispatching many small kernels and --- in most implementations
--- synchronizing with the host to test a termination condition. The relative cost
of launch, synchronization, and Python dispatch against genuine arithmetic is
strongly regime-dependent, so the bottleneck is expected to \emph{move} with scale.
This migration has not been systematically mapped.

\emph{How much is recoverable?} Compilation, fusion, batching, and reduced
precision are standard remedies, but their headroom in NCO is unquantified. Some of
the largest overheads are structurally invisible to the standard remedy:
\texttt{torch.compile} cannot remove a synchronization originating in Python
control flow outside the compiled subgraph.

\emph{Is it worth it?} The Amortized Efficiency Threshold (AET) framing
\citep{aet2026} asks how many instances a neural solver must serve before its
training energy is amortized by its per-instance inference advantage, reporting a
crossover near $1.58\times10^{5}$ instances at $n{=}20$ --- validated, by its
authors' account, at a single scale. Whether the threshold rises or inverts with
scale, batch, and hardware is open.

The stakes are concrete. On TSP-1000, POMO with 8-fold augmentation is reported at
$40.57\%$ gap and $8.5$ minutes, against $3.168\%$ and $1.6$ minutes for a greedy
heatmap-plus-search method \citep{luo2024icam}. The quality--latency--scale
trade-off is thus the dominant design question, and it is most acute in vehicle
routing, where deployment pressure is highest and where the gap between what
learning-based methods promise and what practitioners can operate remains a
recognized open problem \citep{bai2023analytics}.

A substantial efficiency literature exists: StruDiCO reports $3.5\times$ speedup
and $70\%$ memory reduction, and Fast T2T trades one-step generation for
order-of-magnitude gains \citep{wang2025strudico,li2024fastt2t}. Both are
\emph{algorithmic} --- they change the sampler to reduce work. This proposal is
complementary and deliberately \emph{system-level}: it proposes no new solver and
does not compete on quality. It asks what a fixed, representative solver family
costs to run at deployment scale, what fraction of that is recoverable by systems
techniques alone, and how the answer changes when the deployment target changes.

\textbf{Contributions.} (i)~A reproducible characterization of where NCO inference
time is spent and how the hotspot distribution migrates with scale. (ii)~A
quantified account of recoverable headroom, including explicit treatment of the
synchronization bottleneck that compilation cannot address. (iii)~A multi-objective
deployment formulation yielding Pareto-optimal points, hypervolume comparison, and
a decision rule under a quality SLA. (iv)~A generalization study across problem
family and execution backend, with a hardware calibration anchor that makes
cross-backend comparison interpretable rather than merely descriptive.

%======================================================================
\section{Scope and Operational Definitions}
%======================================================================

\subsection{What ``Large-Scale'' Means Here}

In the NCO literature, ``large-scale'' almost always means large \emph{instance
size} $n$. For a deployment question the binding constraint is usually else: how
many instances arrive, how many are served concurrently, and how long the service
runs uninterrupted. We therefore define \textbf{deployment scale} (primary) as the
number of instances $N$, the concurrency $B$, and the sustained operating duration
$\tau$; and \textbf{instance scale} (secondary) as per-instance size $n$, with
training to $n \le 500$ and zero-shot inference generalization to $n = 1000$.

This changes what is measured. Sustained duration becomes a first-class factor, and
it is the factor on which heterogeneous deployment targets are most likely to
diverge, since it exposes thermal and power behavior that short benchmarks cannot
reveal.

\subsection{Objectives}

The objective vector is
$\vect{f}(x) = (\mathrm{gap},\ \mathrm{p95\ latency},\ \mathrm{peak\ memory},\ E_{\mathrm{inst}})$,
evaluated at fixed hardware and batch size. These conflict in the relevant regime:
more samples or deeper search improves gap at the cost of latency and energy;
larger batches improve throughput but raise peak memory.

One scoping decision matters. Only \emph{quality and latency} define the Pareto
front proper. Peak memory is treated as a feasibility constraint --- an operating
point that does not fit is not a candidate --- and energy as a reported attribute
feeding the AET analysis. Four-objective fronts over the sample sizes available
here are dominated by mutually non-dominated points and carry little information.

\subsection{Measured Asymmetries}

Two asymmetries between the deployment targets are not fixable and are stated as
scope boundaries rather than controlled for. \emph{Energy on Apple silicon:} macOS
exposes no per-process GPU energy interface, so energy objectives and the AET
analysis are performed on the CUDA/Linux arm only; the cross-backend comparable
metric set is latency, throughput, and peak memory. \emph{Accelerator memory:}
8\,GB unified versus 11\,GB discrete. Batch size is therefore a grid factor, never
a matched constant, and memory exhaustion is a reportable boundary rather than a
failed run.

%======================================================================
\section{Research Questions and Hypotheses}
%======================================================================

Hypotheses predict a \emph{shape} --- a crossover, a trend, a reversal --- rather
than a direction, since shape claims are the ones most likely to survive as
contributions when specific numbers do not reproduce.

\begin{description}[leftmargin=0pt]

\item[\textbf{RQ1} --- Characterization.] Where is inference time spent, and does
the hotspot distribution migrate with $n$, $B$, and decoding budget?
\textbf{H1.} At small $n$ and $B$ cost is dominated by kernel launch,
host--device synchronization, and Python dispatch; as these grow, cost migrates to
attention and memory bandwidth. Operationally the normalized efficiency $E$
(\S\ref{sec:anchor}) starts well below $1$ and approaches $1$ with scale.
\emph{The crossover scale is itself a primary result}, identifying where the
workload becomes compute-bound.

\item[\textbf{RQ2} --- Optimization headroom.] What is the benefit of compilation,
fusion, batching, and reduced precision, and does it scale with $n$?
\textbf{H2.} Benefit increases with $n$ with clear diminishing returns. More
specifically: gains from \texttt{torch.compile} are \emph{small precisely where H1
predicts overhead dominates}, because the dominant overhead lies outside the
compiled subgraph. Converting the variable-length decode loop to a fixed-length
loop removes the synchronization and enables graph capture. This predicts an
interaction --- compile-only helps little at small scale, loop restructuring helps a
great deal --- and testing it is a principal objective.

\item[\textbf{RQ3} --- Deployment configuration.] Under a quality SLA
($\mathrm{gap} \le \delta$), what is the cost-optimal configuration, and can it be
predicted rather than searched? \textbf{H3.} Latency and memory are predictable
from $(n, B, \text{hardware}, \delta)$ by a fitted cost model whose error is small
relative to variation across configurations, making the decision rule usable
without exhaustive search.

\item[\textbf{RQ4} --- Amortized efficiency.] What is the AET relative to HGS/LKH,
and how does it vary with $n$, $B$, and hardware? \textbf{H4.} AET decreases
substantially as $n$ grows, since the batching advantage of neural inference
amplifies with scale while per-instance training cost is fixed --- a direct
extension of \citep{aet2026} beyond its single validated scale.

\item[\textbf{RQ5} --- Generalization.] Do the observed laws survive a change of
problem family and execution backend, and do optimal configurations transfer?
\textbf{H5.} The \emph{shape} of the trade-off transfers, but optimal operating
points and their ordering do not --- reversals occur. Two mechanisms are
predicted: sustained-drift crossover, where a thermally constrained target leaves
the SLA under long $\tau$ while an actively cooled target does not; and
synchronization-cost asymmetry, where a backend with higher per-op dispatch cost
loses at exactly the small-scale points H1 identifies as overhead-bound.

\end{description}

%======================================================================
\section{Methodology}
%======================================================================

\subsection{System Under Study}

The primary system is RL4CO \citep{berto2023rl4co}, built on PyTorch, TorchRL,
TensorDict, and PyTorch Lightning. At the pinned revision it provides 20 model
families and 24 problem environments. Its value here is threefold: it supplies all
four paradigms of \S\ref{sec:crossmodel} behind one interface; it separates
environment from method, so problem and algorithm vary independently; and its
decode loop is a representative, \emph{unoptimized} autoregressive implementation,
which is what makes it a meaningful object of systems study. The reference workload
for Layers A--C is Attention Model on TSP-100.

\subsection{Layer A: Profiling and Characterization}

Tooling is applied in fixed order, least to most invasive, so observations
cross-validate and measurement perturbation is bounded: \texttt{py-spy} (Python
versus native flame graphs); \texttt{scalene} (CPU/accelerator/memory attribution);
\texttt{perf} (hardware counters, Linux arm); Nsight Systems and \texttt{ncu}
(kernel timeline, occupancy, and the inter-kernel \emph{idle gaps} that reveal
launch and synchronization stalls); and \texttt{pynvml} for power traces.

Each configuration records end-to-end latency (p50/p95/p99), peak accelerator
memory, throughput, and an encode/decode phase split; sustained runs additionally
record a latency drift series. \textbf{Output:} flame graphs, a hotspot ranking, a
hotspot-migration figure, and the $E$ crossover of H1 --- the paper's central
characterization result.

\subsection{Layer B: System-Level Optimization}

Interventions are grouped by priority in Table~\ref{tab:optimizations}, cheapest
and most standard first.

\begin{table*}[t]
\centering
\small
\begin{tabular}{@{}p{0.04\textwidth}p{0.24\textwidth}p{0.64\textwidth}@{}}
\toprule
\textbf{P} & \textbf{Intervention} & \textbf{Description and measurement} \\
\midrule
P0 & \texttt{torch.compile} & Compile the encoder subgraph. Record guard
failures, recompilation count, dynamic-shape penalty. \\
P0 & Batching and padding & Quantify padding waste and the throughput effect of
length bucketing. \\
P0 & Reduced precision (AMP) & Sensitivity of gap and energy to precision.
\emph{CUDA arm only} --- see below. \\
P1 & \textbf{Synchronization elimination in the decode loop} & Replace the
variable-length termination check with a fixed-length loop, using the fact that
decoding length is known a priori for TSP ($n$ steps) and bounded for CVRP. Removes
per-step synchronization, makes shapes static, enables graph capture. Measure
synchronization count, accelerator idle time, end-to-end latency. \\
P2 & ONNX / ONNX Runtime & Static subgraph export; masking and dynamic control
flow complications documented rather than worked around. \\
P3 & TensorRT / INT8 & Conditional. The reference CUDA device is compute
capability 7.5, lacking reduced-precision hardware paths and excluded from current
FlashAttention-2 releases. A negative result is expected and acceptable. \\
\bottomrule
\end{tabular}
\caption{System-level interventions by priority. P1 removes a bottleneck that
compilation cannot reach.}
\label{tab:optimizations}
\end{table*}

\textbf{On precision.} Autocast semantics and operator coverage differ between CUDA
and Metal, and unsupported operations may fall back silently. Reduced precision is
therefore a CUDA-side optimization axis, not a cross-backend comparison condition.
The comparable cross-backend baseline is uniformly FP32
(\texttt{precision=32-true}, \texttt{matmul\_precision=highest}); any
reduced-precision Metal result is reported as exploratory and excluded from
cross-backend claims.

\subsection{Layer C: Multi-Objective Deployment Evaluation}

Following \S2.2, the Pareto front is computed over (gap, p95 latency). An
$\varepsilon$-constraint sweep imposes $\mathrm{gap} \le \delta$ and minimizes
latency, tracing the front per fixed (hardware, batch) slice so that fronts stay
two-dimensional and comparable. Hypervolume compares front quality across
configurations and backends; because HV depends on objective scaling and on the
reference point, both are fixed and reported explicitly --- objectives normalized
per instance scale, gap expressed against best-known, reference point set at the
worst observed value per objective inflated by a fixed margin. Without these
choices stated, HV numbers are neither reproducible nor comparable to future work.

The Amortized Efficiency Threshold is
\begin{equation}
  \mathrm{AET}_E \;=\; \frac{E_{\mathrm{train}}(M)}
  {\max\bigl(E^{\mathrm{b}} - E^{\mathrm{nn}},\ \epsilon\bigr)},
\end{equation}
subject to $\mathrm{gap}(M,\pi) \le \mathrm{gap}(B,\pi) + \delta$, where
$E^{\mathrm{b}}$ is the per-instance energy of the baseline heuristic and
$E^{\mathrm{nn}}$ that of the neural solver under policy $\pi$. Crucially
$E_{\mathrm{train}}(M)$ includes all failed runs and hyperparameter search, not
only the converged model; excluding them understates the numerator by one to two
orders of magnitude and is the most common way AET figures mislead.

\textbf{Output:} a reproducible front with $\ge 5$ points, an HV comparison across
configurations, and an AET surface over $n$, $B$, and hardware. The practical
deliverable is a \emph{deployment decision rule}: given a workload shape and a
quality SLA, which configuration to run, and whether a neural solver is justified
at all against a strong heuristic.

\subsection{Layer D: Generalization}
\label{sec:crossmodel}

\paragraph{Cross-backend (not cross-architecture).} This axis compares CUDA on a
discrete accelerator against Metal/MPS on Apple silicon. \textbf{The comparison must
be described as cross-backend, not cross-architecture.} The two systems differ
simultaneously in instruction set (x86-64 versus ARM64), accelerator API (CUDA
versus Metal), memory system (discrete GDDR6 at $\sim$616\,GB/s versus unified
LPDDR5 shared with the OS), thermal design (actively cooled versus fanless), and
runtime stack (CUDA runtime versus Metal Performance Shaders). No observed
difference can be attributed to processor architecture.

What this axis \emph{can} establish: the existence and location of portability
boundaries, the cost of operator fallback, and whether optimal operating points
transfer. What it \emph{cannot}: any claim about ARM versus x86. A boundary result
is an expected contribution here, framed as a portability finding. A genuine
architecture claim requires a CUDA-capable ARM target, which isolates instruction
set from accelerator API; a Jetson Orin device is requested in \S\ref{sec:resources}
and, if obtained, converts this into a controlled study.

\paragraph{Cross-problem.} Transfer from TSP to CVRP, which adds a capacity
constraint. Training at $n \in \{50,100\}$; evaluation at $n \in \{200,500,1000\}$
is strictly zero-shot, since fine-tuning at test scale would dissolve the question.
The criterion is whether the shape of the trade-off and the ordering of
configurations are preserved under the added constraint.

\paragraph{Cross-model.} Generalization across solver structure is tested
\emph{within} the reference codebase, spanning all four paradigms: autoregressive
(\texttt{am}, \texttt{pomo}, \texttt{symnco}, \texttt{mdam}), non-autoregressive
heatmap (\texttt{nargnn}, \texttt{l2d}), improvement (\texttt{dact},
\texttt{neuopt}, \texttt{eas}), and decomposition (\texttt{glop}). This is
deliberate. Comparing against an independently developed solver introduces a
confound that is easy to overlook: a second codebase brings its own data pipeline,
batching implementation, and instrumentation points, so a measured difference is a
difference between two \emph{implementations}, not two \emph{architectures}, and the
harness itself becomes an uncontrolled variable. Holding the harness fixed isolates
the quantity of interest. Where an external solver is available it serves as an
external validation point --- a law established here is re-tested on it --- which is
a strengthening result if obtained and not a dependency if not.

\subsection{Hardware Calibration Anchor}
\label{sec:anchor}

A raw cross-backend latency ratio is not interpretable on its own. If a workload
takes $4.6\times$ longer on backend $t$ than on reference $r$, that is \emph{equally
consistent} with three different findings: the target hardware is slower, the target
backend is inefficient, or neither machine is doing arithmetic and both are bound
by launch and synchronization overhead. We therefore establish a calibration anchor
--- sustained FP32 matrix multiplication throughput, measured identically on each
target (a 4096$^3$ matmul, synchronized on both sides of the timed region). This is
the most vendor-optimized operation available and bounds achievable compute rate
from above; it is not a predictor of solver performance.

Let $\Pi_r$, $\Pi_t$ be calibrated throughputs. For workload $w$,
\begin{equation}
  H_t = \frac{\Pi_r}{\Pi_t}, \quad
  R_{t,w} = \frac{T_{t,w}}{T_{r,w}}, \quad
  E_{t,w} = \frac{R_{t,w}}{H_t}.
\end{equation}

Table~\ref{tab:E} reads this metric.

\begin{table}[t]
\centering
\footnotesize
\begin{tabular}{@{}l p{0.27\columnwidth}p{0.51\columnwidth}@{}}
\toprule
$E$ & \textbf{Regime} & \textbf{Interpretation} \\
\midrule
$\approx 1$ & Compute-bound & Time ratio tracks silicon; the backend is a faithful,
slower implementation. \\
$\ll 1$ & Overhead-bound & Time is spent outside arithmetic, so latency is similar
despite different hardware. \emph{Evidence for H1.} \\
$\gg 1$ & Backend inefficiency & The target leaves hardware unused --- operator
fallback, per-op synchronization. A portability finding, not a hardware property. \\
\bottomrule
\end{tabular}
\caption{Interpretation of the normalized efficiency metric $E$.}
\label{tab:E}
\end{table}

H1 thereby takes a sharp form: \emph{$E$ should rise from $\ll 1$ at small scale
toward $1$ as $n$ and $B$ grow}, and the crossover scale is reportable. The
interesting question is not the absolute gap between deployments but whether $E$
moves under optimization, whether it moves differently at different scales, and
where it crosses the deployment SLA.

%======================================================================
\section{Experimental Protocol}
%======================================================================

\subsection{Factors and Design}

A full factorial over all factors is neither feasible nor statistically useful for a
characterization study. Layer A runs the full $n \times B$ grid (16 cells) with
decoding fixed, plus a sustained-duration series at anchors; decoding budget and
precision are varied one-at-a-time at a few anchor configurations. Layer B varies
interventions one at a time against the eager baseline on a reduced $n$ grid, so
effects are attributable. Layer C uses the full decoding $\times\ \delta$ product,
which $\varepsilon$-constraint sweeping requires. Table~\ref{tab:factors} lists the
factors and the level at which each is varied.

\begin{table*}[t]
\centering
\small
\begin{tabular}{@{}p{0.14\textwidth}p{0.24\textwidth}p{0.16\textwidth}p{0.38\textwidth}@{}}
\toprule
\textbf{Factor} & \textbf{Levels} & \textbf{Full grid} & \textbf{Notes} \\
\midrule
Problem & TSP, CVRP & Layer D & \\
Instance size $n$ & 50, 100, 500, 1000 & Layer A & train $\le 500$; 1000 zero-shot \\
Batch size $B$ & 1, 8, 32, 128 & Layer A & grid factor, never matched across arms \\
Sustained duration $\tau$ & short, 5, 15, 30 min & Layer A (anchors) & thermal drift \\
Decoding budget & fast / balanced / quality & Layer C & greedy, fixed-$k$, multi-start \\
Execution mode & eager, compiled, restructured & Layer B & OAT against eager \\
Precision & FP32, FP16 (AMP) & Layer B & \textbf{CUDA arm only} \\
Hardware target & CUDA, Metal/MPS & Layer D & (+ Jetson if available) \\
\bottomrule
\end{tabular}
\caption{Experimental factors. ``OAT'': varied one-at-a-time at anchor
configurations.}
\label{tab:factors}
\end{table*}

Decoding budget is evaluated at three interpretable tiers --- \emph{fast} (greedy or
single sample), \emph{balanced} (fixed sample count), \emph{quality} (fixed time or
search depth) --- yielding a decision curve rather than a point estimate.

\subsection{Instance Sets and Determinism}

Instance sets are generated once with fixed seeds using the reference codebase's
public generators, serialized to disk, and loaded by every subsequent run. They are
\emph{never} regenerated per run and \emph{never} regenerated per backend: a
cross-backend comparison on independently sampled instances compares different
problems. External validity uses a TSPLIB subset and CVRPLIB Set-X against official
optima or best-known solutions. Datasets, checkpoints, and logs are written outside
any cloud-synchronized directory --- background synchronization during a timed run
perturbs wall-clock measurements invisibly.

\subsection{Measurement and Reproducibility}

Fixed warm-up ($\ge 3$ iterations) discards first-call compilation and clock
ramp-up; $\ge 30$ repetitions per configuration give mean, standard deviation, and
95\% confidence intervals. Comparisons are \emph{paired} on the same instance set
and seed across arms, using common random numbers to reduce variance. Instance-level
and run-level variability are reported \emph{separately}: thirty repetitions of the
same 100 instances are not 300 independent samples for instance-level effects. The
accelerator is synchronized before reading the clock --- CUDA and MPS both queue
asynchronously, so without synchronization one measures enqueue, not compute. Energy
is integrated as $\int P\,dt$, not averaged power. Clocks are locked where possible;
otherwise the full dynamic frequency trace is recorded. For sustained runs the
calibration anchor is re-measured before and after, since a thermally constrained
target throttles and the throttled throughput would otherwise be silently attributed
to the backend.

Reproducibility is a deliverable, not an afterthought. The codebase is pinned by
revision hash and held read-only; environment problems are resolved in the
environment, never by editing source, since a measurement must correspond to a
specific unmodified revision. Each machine produces a machine-generated environment
record (interpreter and package versions, revision, hardware capability, calibration
anchor), and every table and figure is regenerated by script from raw measurements
--- no number is transcribed by hand.

\subsection{Metrics}

\textbf{Optimality gap}: $(c_{\text{model}} - c_{\text{ref}})/c_{\text{ref}}$,
reported as mean, quantiles, and failure count. Where gap saturates at zero (small
$n$, strong solvers) we report the fraction of instances solved to best-known rather
than a degenerate zero. \textbf{Latency}: end-to-end p50/p95/p99, split into encode
and decode. \textbf{Throughput}: instances/s, and samples/s during sampling, at fixed
$n$ so large instances are not confounded with low counts. \textbf{Peak memory}:
maximum accelerator memory allocated, reset per configuration. \textbf{Energy}:
J/instance and J per unit gap improvement, CUDA/Linux arm only. \textbf{Training
cost}: wall-clock per epoch and total cost to reach target gap, with optimizer,
batch size, and parallel environment count reported. \textbf{Deployment cost}: AET
in instances, with $\varepsilon$ and grid carbon intensity reported.

Reference solutions come from LKH-3 for TSP \citep{helsgaun2017lkh3} and PyVRP for
CVRP \citep{wouda2024pyvrp}, the latter also being the baseline used by the AET
framework, preserving comparability with that line of work. OR-Tools
\citep{perron2023ortools} is used for independent feasibility checking. An exact
solver (Gurobi, academic license held) is viable for certified-optimal references
only at small $n$ and for tightly constrained CVRP; it is \emph{not} a viable gap
denominator for TSP at $n \ge 500$, where it cannot close the gap in practical time.

%======================================================================
\section{Resources, Timeline, and Risks}
\label{sec:resources}
%======================================================================

\subsection{Hardware, Limits, and Requests}

The available hardware is an RTX 2080 Ti (11\,GB, Turing, sm\_75) and a MacBook Air
M2 (8\,GB unified, fanless). Verified properties matter here: the 2080 Ti executes
native compiled code with no PTX JIT fallback, but 11\,GB constrains $n{=}500$
training with multi-start decoding, it has no reduced-precision hardware path, and
it is excluded from current FlashAttention-2 releases. The M2 has partial operator
coverage with silent CPU fallback, cannot represent edge ARM deployment generally,
exposes no per-process accelerator energy, and throttles under sustained load.

Note the asymmetry this creates. Jointly the two devices suit Layers A, B, and the
cross-backend part of D well --- and the thermal contrast between an actively cooled
and a fanless device is, unusually, an \emph{asset} for the sustained-duration
factor. They are not suited to large-scale training, symmetric energy measurement,
or any instruction-set architecture claim.

Three requests follow, in order of necessity. \textbf{An Ubuntu installation or
dual-boot on the CUDA machine} --- RQ4 requires energy measurement via RAPL, which
has no supported interface on Windows; this is a hard prerequisite for RQ4, not a
convenience. \textbf{An NVIDIA Jetson Orin or ARM server} --- the only way to
isolate instruction set from accelerator API; without it the architecture axis
remains a backend comparison. \textbf{A GPU with $\ge 24$\,GB or A100 node access}
--- unlocks $n{=}500$ training and the larger batch sizes. Helpful: PC memory to
32\,GB for $n{=}1000$ dataset generation.

\subsection{Timeline}

The study is designed so that a defensible contribution exists at several exit
points, not only at the end. The \textbf{minimum viable result}, achievable on
current hardware, is: reproducible profiling and benchmarking code; a characterized
hotspot migration with the $E$ crossover; measured optimization headroom; $\ge 5$
Pareto-optimal points with hypervolume comparison; and an honest record of negative
results. The plan in Table~\ref{tab:timeline} is sequenced to protect that result.

\begin{table*}[!tb]
\centering
\footnotesize
\begin{tabular}{@{}p{0.08\textwidth}p{0.17\textwidth}p{0.40\textwidth}p{0.25\textwidth}@{}}
\toprule
\textbf{Month} & \textbf{Phase} & \textbf{Deliverable} & \textbf{Go/No-go} \\
\midrule
M1 & Environment freeze, baseline & Pinned environment records per machine; RL4CO
TSP-50/100 running; LKH-3, PyVRP, OR-Tools references; calibration anchor on both
targets. & Reproduced gap matches published magnitude; anchors reproducible to a
few percent. \\
M2 & Characterization (A) & Flame graphs, hotspot ranking, $n \times B$ grid,
sustained-duration series, hotspot-migration figure, $E$ crossover. & Migration
visible; $E$ trend monotone as predicted. \\
M3 & Optimization I (B, P0) & Compile, batching, precision ablations against eager.
& Attribution clean: one intervention at a time. \\
M4 & Optimization II and multi-objective (B P1, C) & Decode-loop synchronization
elimination with graph capture; $\varepsilon$-constraint sweep; Pareto front and
hypervolume; decision rule. & $\ge 5$ Pareto points; \emph{minimum viable result
complete}. \\
M5 & Deployment cost and cross-model (RQ4) & AET computation and sensitivity;
cross-paradigm generalization; external validation if available. & AET surface over
$n$, $B$, hardware; Linux environment in place. \\
M6 & Cross-backend, cross-problem, writing (D) & TSP to CVRP transfer; backend
comparison with reversal analysis; Jetson results if available; draft. & All tables
regenerable from scripts end-to-end. \\
\bottomrule
\end{tabular}
\caption{Six-month plan; a complete result is targeted at M4.}
\label{tab:timeline}
\end{table*}

\subsection{Risks}

Each risk carries a pre-committed mitigation and an early-warning signal, so the
response is decided before the signal appears rather than after it. The signal
column exists so that a risk is detected when it first becomes true, not when it
becomes fatal. Table~\ref{tab:risks} lists them.

\begin{table*}[!tb]
\centering
\footnotesize
\begin{tabular}{@{}p{0.20\textwidth}p{0.18\textwidth}p{0.54\textwidth}@{}}
\toprule
\textbf{Risk} & \textbf{Signal} & \textbf{Mitigation} \\
\midrule
Memory insufficient & OOM at $n{=}500$ training & Restrict $n{=}1000$ to inference;
gradient checkpointing; FP16; smaller batch. \emph{Do not modify model architecture
to mask a systems problem.} \\
\texttt{torch.compile} yields little & Speedup $< 5\%$ & Expected if H1 holds.
Pivot weight to decode-loop restructuring (P1) and the AET cost model, which are the
primary contributions regardless. The negative result becomes evidence \emph{for}
H1. \\
Metal not comparable & Substantial operator fallback & Report as a portability
boundary. Silent-fallback characterization is a legitimate result and does not
invalidate the CUDA-side findings. \\
RQ4 blocked & No Linux environment by M5 & Defer RQ4; Layers A--C do not depend on
it. \\
Scope overrun & Layer B not converged by end of M3 & Cut P2 and P3; drop the Jetson
arm; retain the two-backend comparison and the complete Pareto analysis. The minimum
viable result is intact. \\
No external solver available & Not obtained by week 4 & The cross-paradigm axis
within the reference codebase is self-sufficient and, as argued in
\S\ref{sec:crossmodel}, methodologically cleaner. This is the designed fallback, not
a degradation. \\
\bottomrule
\end{tabular}
\caption{Risks, signals, and pre-committed mitigations.}
\label{tab:risks}
\end{table*}

\clearpage
\bibliographystyle{abbrvnat}
\bibliography{references}

\end{document}
